%%
%% fall-lecture09.tex
%% 
%% Made by Alex Nelson <pqnelson@gmail.com>
%% Login   <alex@lisp>
%% 
%% Started on  2026-10-17T11:01:54-0700
%% Last update 2026-10-17T11:01:54-0700
%% 

\lecture{}

\begin{node}
We will discuss compactness in the space of continuous [real-valued]
functions $C(X)$.  This is important for three reasons
\begin{enumerate}
\item It illustrates the Hausdorff theorem
\item The space is important in its own right (e.g., linear
  functionals on such spaces give another way to do measure theory)
\item This will be important next quarter.
\end{enumerate}
\end{node}

\begin{notation}
Let $X$ be a set, let $(Y,\rho)$ be a metric space. We denote by $M(X,Y)$ the
set of bounded maps from $X$ to $Y$ (i.e., the $f$ such that $f(X)$ is
bounded in $Y$).

We can define the function $d\colon M(X,Y)\times M(X,Y)\to[0,\infty)$
by
\begin{equation}
d(f,g) := \sup_{x\in X}\rho(f(x),g(x)).
\end{equation}
We claim $d$ is a metric, and it induces the topology of uniform
convergence on $M(X,Y)$. (Exercise: verify $d$ really is a metric.)
\end{notation}

\begin{remark}
We can equip $M(X,Y)$ with the \define{Topology of Uniform Convergence} which 
is generated by sets of the form
\begin{equation}
A_{n,f}\{g\in M(X,Y)\mid\sup_{x\in X}\rho(g(x),f(x))<\frac{1}{n}\}
\end{equation}
for any $n\in\NN$ and $f\in M(X,Y)$. 
\end{remark}

\begin{proposition}
If $(Y,\rho)$ is a complete metric space, then the metric $d$ is complete.
\end{proposition}

\begin{example}
We see $M(X,\RR)$ is complete since $\RR$ is complete.
\end{example}

\begin{node}
We denote by $C_{b}(X)$ the subspace of continuous bounded
[real-valued] functions on $X$. It is a subspace of $C(X)$.
\end{node}

\begin{theorem}[Weierstrass uniform convergence]
The subspace of bounded continuous functions $C_{b}(X)$ on a
topological space $X$ may be embedded into $C_{b}(X)\into M(X,\RR)$
(which is obvious) and the $C_{b}(X)$ is a closed subset of $M(X,\RR)$.

In particular, if $X$ is compact, then [continuous functions are
  automatically bounded so] $C_{b}(X)=C(X)$.
\end{theorem}

``Uniform limit of continuous functions is continuous'' is the slogan.

\begin{theorem}[Arzela--Ascoli]
Let $X$ be a compact topological space. Let $\mathcal{F}\subset C(X)$
be a family of continuous maps.
Then $\mathcal{F}$ is pre-compact [i.e., its closure
  $\closure{\mathcal{F}}$ in $C(X)$ is compact] if and only if both
(1)~$\mathcal{F}$ is equicontinuous and (2)~$\mathcal{F}$ is pointwise-bounded.
[By the Hausdorff theorem, precompactness is equivalent to ``totally bounded''.]
\end{theorem}

\begin{remark}
Recall, $\mathcal{F}$ is \define{Equicontinuous} if it is
equicontinuous at every point $x_{0}$---i.e., for any $\varepsilon>0$
there exists a neighborhood $U\ni x_{0}$ in $X$ such that if
$f\in\mathcal{F}$ and [usual notion of continuity at $x_{0}$ holds for
  $f$, i.e.] $z\in U$ then $|f(z)-f(x_{0})|<\varepsilon$.
Observe this means that we use the same $\varepsilon$-$\delta$
parameters for all $f\in\mathcal{F}$.
\end{remark}

\begin{remark}
Regarding $\mathcal{F}$ being pointwise bounded, this means: for every
point $x\in X$, $\{f(x)\mid f\in\mathcal{F}\}$ is bounded in $\RR$.
\end{remark}

\begin{example}
Let $X=[0,1]$ be given the subspace topology (considered as a subspace
of $\RR$ with standard topology). 

Let $\mathcal{F}_{1}=\{x^{n}\mid n\in\NN\}$ be monomial functions. Let
$\mathcal{F}_{2}=\{f(x)=n\mid n\in\NN\}$ be constant functions.

Then $\mathcal{F}_{1}$ is not equicontinuous (problems happen at
$x=1$).
And $\mathcal{F}_{2}$ is not pointwise-bounded.

So there is no convergence subsequence of \emph{either} family.
\end{example}

\begin{proof}[Proof (Arzela--Ascoli theorem)]
\forwardproof\ Assume $\mathcal{F}$ is precompact. Then $\mathcal{F}$
is totally bounded. Then there exists $f_{1}$, \dots, $f_{N}\in\mathcal{F}$
such that every $f\in\mathcal{F}$ is a distance $0<\varepsilon\leq1$
from one of these $d(f,f_{j})\leq1$ for some
$j\in\{1,\dots,N\}$. Since each $f_{j}$ is continuous, this means
$f_{1}(X)$, \dots, $f_{N}(X)$ are all compact. Then there exists some
radius $R>0$ such that
\begin{equation}
\bigcup^{N}_{j=1}f_{j}(X)\subset(-R,R)
\end{equation}
the images sit inside an open ball of radius $R$ in $\RR$. This gives
us uniform boundedness (which implies pointwise-boundedness).

Now, for equicontinuity, let $x_{0}\in X$ be a point. Let
$\varepsilon>0$ be arbitrary. We have just proven total boundedness,
so there exists functions $f_{1}$, \dots, $f_{N(\varepsilon)}\in\mathcal{F}$
making the set $\{f_{1},\dots,f_{N(\varepsilon)}\}$
$\varepsilon$-dense in $\mathcal{F}$. Since each function $f_{j}$ is
continuous (for each $j=1,\dots,N(\varepsilon)$), then for each $j$
there is a neighborhood $U_{j}(x_{0})\ni x_{0}$ such that every $x\in U_{j}(x_{0})$
satisfies $|f_{j}(x)-f_{j}(x_{0})|<\varepsilon$. Then we can find a
neighborhood $U$ containing each $U_{j}$. Then for all
$f\in\mathcal{F}$, we have
\begin{equation}
|f(x)-f(x_{0})|\leq|f(x)-f_{j}(x)|+|f_{j}(x)-f_{j}(x_{0})|+|f_{j}(x_{0})-f(x_{0})|<3\varepsilon
\end{equation}
where we approximate $f(x)$ by $f_{j}(x)$ with error
$|f(x)-f_{j}(x)|<\varepsilon$. Similarly, we approximate $f(x_{0})$ by
$f_{j}(x_{0})$ with error $|f_{j}(x_{0})-f(x_{0})|<\varepsilon$, and
we know $|f_{j}(x)-f_{j}(x_{0})|<\varepsilon$. We see that this
suffices to prove equicontinuity.

\backwardproof\ Assume $\mathcal{F}$ is equicontinuous and
pointwise-bounded. We want to show $\Clos_{C(X)}(\mathcal{F})$ is
compact. Since (by Weierstrass theorem),
\begin{equation}
\Clos_{C(X)}(\mathcal{F})=\Clos_{C_{b}(X)}(\mathcal{F})=\Clos_{M(X,\RR)}(\mathcal{F}),
\end{equation}
we just need to show that $\mathcal{F}$ is compact, which is
equivalent to showing $\mathcal{F}$ is totally-bounded in $M(X,\RR)$.
[We do this because it is easier to approximate by bounded functions.]

Let $\varepsilon > 0$. Using comapctness of $X$, we can represent
\begin{equation}
X=\bigsqcup_{j}X_{j}
\end{equation}
as a finite disjoint union of sets $X_{j}\subset U(x_{j})$ where
$U(x_{j})$ are neighborhoods taken by definition of equicontinuity. We
denote
\begin{equation}
L:=\Lin\{\mathbf{1}_{X_{j}}\}
\end{equation}
the $\RR$-linear combination functions of the characteristic functions
$\mathbf{1}_{X_{j}}$ of the subsets $X_{j}$, since
$\mathbf{1}_{X_{j}}\in M(X,\RR)$. So this means that $L\iso\RR^{N}$ as
vector spaces. So $L\subset M(X,\RR)$. If $f\in\mathcal{F}$, then the
approximating function
\begin{equation}
\widetilde{f}=\sum^{N}_{j=1}f(x_{j})\mathbf{1}_{X_{j}}
\end{equation}
belongs to $\widetilde{f}\in L$, and $d(f,\widetilde{f})<\varepsilon$
by equicontinuity. This gives us an estimate. Therefore $\mathcal{F}$
is a distance $\varepsilon$ from an open ball in $L$, and this ball
has radius $\|\mathcal{F}\|$. Hence the claim (of precompactness).
\end{proof}

\begin{remark}
In the textbook, several versions of this theorem are given. One
generalization is to locally compact Hausdorff spaces (called ``LCH
spaces'' in the book). There are also LCH generalizations of many
useful theorems.
\end{remark}

\begin{definition}
We say a topological space $X$ is \define{Locally Compact} if every
point has a compact neighborhood.
\end{definition}

\begin{example}
We see $\RR^{n}$ is locally compact, for any $n\in\NN$.
\end{example}

\begin{example}
Open sets of LCH spaces are locally compact.
\end{example}

\subsection{1-Point Compactification}

\begin{node}
Let $(X,\topology)$ be a topological space. Then we can add one more point
[which is not already in $X$] to it
and call this new point ``infinity'' (and denote it $\infty$). Then
the set
\begin{equation}
X_{\infty}:=X\sqcup\{\infty\}
\end{equation}
may be given the following topology:
\begin{equation}
\topology_{\infty}=\topology\cup\{X\setminus K\mid K\subset X\mbox{ is compact}\}.
\end{equation}
Open sets in $X_{\infty}$ are either open sets in $X$ \emph{or}
complements of compact subsets of $X$.
\end{node}

\begin{example}
We see that $\RR_{\infty}\sim S^{1}$ is the unit circle, and
$\RR^{n}_{\infty}\sim S^{n}$.
\end{example}

\begin{claim}
If $(X,\topology)$ is LCH, then $(X_{\infty},\topology_{\infty})$ is Hausdorff and the relative
topology $(\topology_{\infty})_{X}$ of $X$ in $X_{\infty}$ coincides with
$\topology=(\topology_{\infty})_{X}$. 
\end{claim}

\begin{remark}
There are other compactifications, e.g., the Stone-\v{C}ech
compactification. More generally~\cite[p.144]{folland1999real}, a compactification of $X$ will be a
pair $(Y,\varphi)$ where $Y$ is a compact Hausdorff(!) space and
$\varphi$ is a homeomorphism of $X$ onto a dense subset of $Y$.
\end{remark}